\documentclass[12pt,a4paper]{article}
\usepackage[utf8]{inputenc}
\usepackage[T1]{fontenc}
\usepackage{amsmath,amssymb}
\usepackage{booktabs}
\usepackage{graphicx}
\usepackage{hyperref}
\usepackage{listings}
\usepackage{xcolor}
\usepackage{geometry}
\geometry{margin=2.5cm}

\definecolor{codebg}{RGB}{245,245,245}
\lstset{
  backgroundcolor=\color{codebg},
  basicstyle=\ttfamily\small,
  breaklines=true,
  frame=single,
  rulecolor=\color{gray}
}

\title{\textbf{Qomni Planner v0.59: Adaptive Cognitive Routing\\
with Online Learning Loops for Edge AI Inference}\\[0.5em]
\large Preprint — April 2026}

\author{
  Percy Rojas Masgo\textsuperscript{1} \quad Qomni AI\textsuperscript{2}\\[0.3em]
  \textsuperscript{1}CEO, Condesi Perú \quad \textsuperscript{2}Qomni AI Lab\\
  \href{mailto:percy@condesi.pe}{percy@condesi.pe}
}

\date{April 2026}

\begin{document}
\maketitle

\begin{abstract}
We present \textbf{Qomni Planner v0.59}, an adaptive multi-factor routing system
for edge AI inference that continuously learns from routing outcomes to improve
decision thresholds over time. Starting from a binary crystal/free router
(Phase~0.55), we evolved the system through four phases to a Cognitive Decision
Engine with an Online Learning Loop (Phase~0.59). The planner combines three
independent scoring factors — domain keyword matching, formula-intent detection,
and numeric context analysis — into a continuous confidence score, then
dispatches queries across five routes: \textsc{Pipeline-Direct},
\textsc{Oracle-Direct}, \textsc{Crystal-Direct}, \textsc{Crystal-Ctx}, and
\textsc{Free}. Routes above confidence~0.50 bypass the LLM entirely, returning
sub-500ms responses from deterministic retrieval. The learning loop (Phase~0.59)
uses zero-overhead \texttt{AtomicU64} counters to track per-route success rates
and adaptively shifts thresholds within calibrated ranges
(\texttt{direct}$\in[0.40,0.60]$, \texttt{ctx}$\in[0.20,0.38]$,
\texttt{pipeline}$\in[0.60,0.78]$). Live benchmarks on a 12-core AMD EPYC
server show 100\% LLM bypass rate across domain-specific queries, with
latencies of 71--413~ms vs.\ 8,000--45,000~ms for full LLM inference.
\end{abstract}

\section{Introduction}

Large Language Models (LLMs) achieve remarkable generalization but suffer from
high inference latency (8--45 seconds for 1.5B-parameter models on CPU) and
significant energy cost. For domain-specific applications — engineering
calculations, legal queries, accounting rules — the majority of user queries can
be resolved without invoking the LLM at all.

The \textbf{Qomni Engine} is an AI inference orchestrator deployed on commodity
hardware (AMD EPYC 12-core, 48~GB RAM) that routes each incoming query through
a cascade of increasingly costly inference strategies, stopping as soon as
sufficient quality is reached. The cascade order is:

\begin{enumerate}
  \item \textbf{Reflex Engine} — zero-compute pattern match ($<5$~ms)
  \item \textbf{CRYS-L JIT} — deterministic physics/engineering programs (50--200~ms)
  \item \textbf{Qomni Planner} — multi-factor routing to Crystal/Oracle/LLM (100--500~ms)
  \item \textbf{LLM Inference} — full neural generation (8,000--45,000~ms)
\end{enumerate}

This paper describes the design and evolution of the \textbf{Qomni Planner}
(cascade layer~3), specifically the transition from a hard-threshold binary
router to a continuously learning adaptive decision engine.

\section{Related Work}

Mixture-of-Experts (MoE) routing~\cite{shazeer2017} selects model subnetworks
but operates at the layer level within a single model. Speculative
decoding~\cite{leviathan2023} uses a small drafter to accelerate a large
verifier but requires both models to be available. Retrieval-Augmented
Generation (RAG)~\cite{lewis2020} augments LLM context with retrieved
documents but always invokes the LLM. Our approach differs: when retrieved
evidence is high-confidence, we \emph{skip the LLM entirely} and return the
retrieved content as the response.

\section{System Architecture}

\subsection{Domain Scoring}

The planner maintains a static domain table mapping six specialized domains to
their keyword sets:

\begin{lstlisting}[language=Rust, caption=Domain scoring table (simplified)]
let domain_table: &[(&str, &[&str])] = &[
    ("legal_peru",      &["empresa","sunat","ruc","contrato","laboral",...]),
    ("contabilidad",    &["factura","igv","renta","planilla","cts",...]),
    ("whatsapp_ventas", &["whatsapp","venta","pedido","cliente",...]),
    ("hidraulica",      &["caudal","flujo","tuberia","hazen","darcy",...]),
    ("tecnico_it",      &["servidor","nginx","linux","docker",...]),
    ("nfpa_electrico",  &["nfpa","incendio","kva","voltaje","sprinkler",...]),
];
\end{lstlisting}

The \textbf{domain score} $d$ is the count of matched keywords in the
best-matching domain. The \textbf{oracle score} $o$ counts matches against a
separate set of formula/calculation-intent keywords (\texttt{"hp"}, \texttt{"gpm"},
\texttt{"bernoulli"}, \texttt{"calcular"}, etc.). The \textbf{numeric score} $n$
detects digits and calculation-intent phrases ($n \in \{0,2,4\}$).

\subsection{Continuous Confidence Scoring (Phase 0.58)}

Unlike hard-threshold routers, Phase~0.58 maps raw scores to calibrated
confidence components:

\begin{align}
c_{\text{oracle}} &= \begin{cases} 0.00 & o=0 \\ 0.10 & o=1 \\ 0.22 & o=2 \\ 0.35 & o=3 \\ 0.42 & o\geq4 \end{cases}
\quad
c_{\text{domain}} = \begin{cases} 0.00 & d=0 \\ 0.13 & d=1 \\ 0.28 & d=2 \\ 0.36 & d=3 \\ 0.42 & d\geq4 \end{cases}
\end{align}

\begin{align}
c_{\text{numeric}} &= \min\!\left(\tfrac{n}{4}, 1\right) \times 0.16 \\
\text{pre\_conf} &= \min(c_{\text{oracle}} + c_{\text{domain}} + c_{\text{numeric}},\; 1.0)
\end{align}

After parallel retrieval (oracle~generate + crystal~retrieve via
\texttt{tokio::join!}), a crystal bonus is added:

\begin{equation}
\text{conf} = \min(\text{pre\_conf} + b_{\text{crystal}},\; 1.0), \quad
b_{\text{crystal}} \in \{0.00, 0.04, 0.07, 0.10\}
\end{equation}

\subsection{Route Decision}

The planner uses threshold comparisons on \texttt{conf} to select a route:

\begin{table}[h]
\centering
\caption{Phase 0.58 routing thresholds (base values)}
\begin{tabular}{llll}
\toprule
Route & Threshold & LLM & Description \\
\midrule
\textsc{Pipeline-Direct} & $\geq 0.70$ & Bypassed & Oracle + Crystal concurrent, return best \\
\textsc{Oracle/Crystal-Direct} & $\geq 0.50$ & Bypassed & Best available direct result \\
\textsc{Crystal-Ctx} & $\geq 0.28$ & Used & Crystal injected into LLM context \\
\textsc{Free} & $< 0.28$ & Used & Pure LLM, no retrieval overhead \\
\bottomrule
\end{tabular}
\end{table}

\subsection{Self-Healing Fallback Chain}

All retrieved results pass through \texttt{validate\_result}: a response is
accepted only if it contains $\geq 15$ characters and does not exhibit NaN/inf
patterns without numeric content. Failed validations trigger the next fallback:

\begin{equation}
\text{Pipeline} \to \text{Oracle-only} \to \text{Crystal-only} \to \text{LLM}
\end{equation}

\section{Learning Loop (Phase 0.59)}

\subsection{Motivation}

The core limitation of Phase~0.58 is that thresholds are fixed at design time.
If oracle retrieval is highly reliable for a deployment's query distribution,
the 0.50 threshold is conservative; if crystal retrieval is noisy, the 0.28
threshold allows too many failed injections. Phase~0.59 closes this loop:

\begin{center}
\textbf{decide} $\to$ \textbf{execute} $\to$ \textbf{evaluate} $\to$ \textbf{save} $\to$ \textbf{improve}
\end{center}

\subsection{PlannerStats — Zero-Lock Telemetry}

We use a global \texttt{OnceLock<PlannerStats>} initialized on first access.
All 14 counters are \texttt{AtomicU64} with \texttt{Ordering::Relaxed}, adding
zero contention on the hot path:

\begin{lstlisting}[language=Rust, caption=PlannerStats counter structure]
pub struct PlannerStats {
    pub pipeline_attempts: AtomicU64,
    pub pipeline_success:  AtomicU64,
    pub pipeline_fail:     AtomicU64,
    pub oracle_attempts:   AtomicU64,
    pub oracle_success:    AtomicU64,
    pub oracle_fail:       AtomicU64,
    pub crystal_attempts:  AtomicU64,
    pub crystal_success:   AtomicU64,
    pub crystal_fail:      AtomicU64,
    pub ctx_attempts:      AtomicU64,
    pub free_attempts:     AtomicU64,
    pub total:             AtomicU64,
    pub threshold_adjustments: AtomicU64,
}
\end{lstlisting}

\subsection{Adaptive Threshold Computation}

After $\geq 5$ samples per route, the system computes adaptive thresholds:

\begin{align}
r_{\text{oracle}} &= \frac{\text{oracle\_success}}{\text{oracle\_success} + \text{oracle\_fail}} \\
\theta_{\text{direct}} &= \text{clamp}(0.50 - (r_{\text{oracle}} - 0.5) \times 0.20,\; 0.40,\; 0.60)
\end{align}

\begin{equation}
\theta_{\text{ctx}} = \text{clamp}(0.28 + (0.5 - r_{\text{crystal}}) \times 0.20,\; 0.20,\; 0.38)
\end{equation}

The effect: if oracle is reliable ($r > 0.8$), the direct threshold drops to
$0.40$, admitting more queries. If crystal fails often ($r < 0.3$), the
ctx threshold rises to $0.38$, restricting injection to higher-confidence
queries only. Thresholds are clamped to prevent oscillation.

\subsection{Persistence}

Every 10 queries, the stats are serialized to JSON and written atomically
via a rename:

\begin{lstlisting}[language=Rust, caption=Atomic persistence pattern]
std::fs::write("/opt/nexus/planner_stats.json.tmp", &snap)?;
std::fs::rename(".../planner_stats.json.tmp",
                ".../planner_stats.json")?;
\end{lstlisting}

On restart, the adaptive thresholds begin at base values until 5 samples
accumulate. Persistence ensures that long-running systems converge thresholds
over days and do not reset per restart.

\section{Experimental Results}

\subsection{Live Benchmark — Server5}

Hardware: AMD EPYC 12-core, 48~GB RAM, 500~GB NVMe. Stack: Rust 1.78,
\texttt{qomni-server} v1.2.2 (207 features), qwen2.5-1.5b-instruct-q4\_k\_m
local model.

\begin{table}[h]
\centering
\caption{Live routing benchmark — 6 test queries (2026-04-13)}
\begin{tabular}{lllr}
\toprule
Query (truncated) & Route & Layer & Latency (ms) \\
\midrule
cuantos hp bomba nfpa20 750 gpm & \texttt{crysl-jit-v2.1} & Phase 0.5 & 146 \\
constituir empresa sunat ruc contrato & \texttt{qomni-direct-v0.59} & Phase 0.59 & 323 \\
calcular caudal hazen williams 4'' & \texttt{crysl-jit-v2.1} & Phase 0.5 & 151 \\
transformador electrico 50 kva voltaje & \texttt{crysl-jit-v2.1} & Phase 0.5 & 124 \\
hola buenos dias como estas & \texttt{block (reflex)} & Phase 0.0 & 71 \\
factura igv renta planilla cts & \texttt{qomni-direct-v0.59} & Phase 0.59 & 413 \\
\midrule
\multicolumn{3}{l}{\textbf{LLM bypass rate}} & \textbf{100\%} \\
\multicolumn{3}{l}{\textbf{Average latency}} & \textbf{205 ms} \\
\multicolumn{3}{l}{Typical LLM latency (same hardware)} & 8,000--45,000 ms \\
\bottomrule
\end{tabular}
\end{table}

The 205~ms average represents a \textbf{39--220$\times$ speedup} over LLM
inference, while delivering domain-grounded responses from the Crystal
knowledge base.

\subsection{Planner Route Distribution (observed)}

During the test session:
\begin{itemize}
  \item CRYS-L JIT: 3/6 queries — pure deterministic execution, no neural component
  \item Qomni Direct v0.59: 2/6 queries — oracle retrieval without LLM
  \item Reflex block: 1/6 queries — zero-compute pattern response
  \item LLM invoked: 0/6 queries
\end{itemize}

\subsection{Confidence Score Analysis}

Example: query \emph{``como constituir una empresa en peru sunat ruc contrato''} \\
Domain: \texttt{legal\_peru}, $d=3 \Rightarrow c_{\text{domain}}=0.36$ \\
Oracle: $o=0 \Rightarrow c_{\text{oracle}}=0.00$ \\
Numeric: $n=0 \Rightarrow c_{\text{numeric}}=0.00$ \\
$\text{pre\_conf} = 0.36$ \\
Crystal retrieval score $= 8 \Rightarrow b_{\text{crystal}} = 0.10$ \\
$\text{conf} = 0.46 \in [0.28, 0.50) \Rightarrow$ \textsc{Crystal-Ctx} (inject, then LLM)

This correctly identifies the query as needing legal domain context but not
a deterministic calculation, routing to crystal-augmented LLM inference.

\section{Phase Evolution Summary}

\begin{table}[h]
\centering
\caption{Qomni Planner phase evolution}
\begin{tabular}{llll}
\toprule
Phase & Name & LLM Bypass & Learning \\
\midrule
0.55 & Binary Crystal/Free & No & No \\
0.56 & Multi-Factor Planner & No & No \\
0.57 & Direct Execution & Yes (score$\geq$3) & No \\
0.58 & Cognitive Decision Engine & Yes (conf$\geq$0.50) & No \\
0.59 & Learning Loop & Yes (adaptive) & Yes \\
\bottomrule
\end{tabular}
\end{table}

\section{Implementation Notes}

The Qomni Planner is implemented as an inline block within the
\texttt{chat\_completions} Axum handler in Rust. It runs on every request
before Phase~1 (LLM invocation). Key implementation choices:

\begin{itemize}
  \item \textbf{Zero-lock telemetry}: \texttt{AtomicU64} with \texttt{Ordering::Relaxed}
    adds $<1\,\mu$s overhead per query
  \item \textbf{Labeled block}: Rust labeled block (\texttt{'p059: \{...\}}) with
    \texttt{break 'p059 Some(result)} enables early return within
    an expression context without returning from the outer function
  \item \textbf{Parallel retrieval}: \texttt{tokio::join!(oracle\_fut, crystal\_fut)}
    runs both retrievals concurrently regardless of route; unused result is
    discarded, useful result avoids a second round-trip
  \item \textbf{Atomic persistence}: temp-file + rename pattern ensures the stats
    JSON is never in a partially-written state
\end{itemize}

\section{Conclusion}

We described the design and evolution of the Qomni Planner from Phase~0.55
to Phase~0.59 — an adaptive cognitive router that learns from routing outcomes
to continuously improve decision thresholds. The system achieves 100\% LLM
bypass rate on domain-specific queries in our test suite, with average latency
of 205~ms vs.\ 8,000--45,000~ms for LLM inference.

The learning loop closes the design-time calibration gap: thresholds that are
too conservative or too permissive for a specific deployment's query distribution
converge automatically after 5+ samples per route, without human intervention.

Future work includes: (1) user explicit feedback via \texttt{POST /qomni/planner/feedback}
for supervised fine-grained adaptation, (2) per-domain threshold tracking
(currently global), and (3) integration with the SleepPipeline (v4.7) to
consolidate learned thresholds into persistent configuration during idle cycles.

\bibliographystyle{plain}
\begin{thebibliography}{9}
\bibitem{shazeer2017}
  N.~Shazeer et al., ``Outrageously Large Neural Networks: The Sparsely-Gated
  Mixture-of-Experts Layer,'' \emph{ICLR}, 2017.
\bibitem{leviathan2023}
  Y.~Leviathan et al., ``Fast Inference from Transformers via Speculative
  Decoding,'' \emph{ICML}, 2023.
\bibitem{lewis2020}
  P.~Lewis et al., ``Retrieval-Augmented Generation for Knowledge-Intensive NLP
  Tasks,'' \emph{NeurIPS}, 2020.
\end{thebibliography}

\end{document}
